{\color{red}
\section{Open questions for the author (NOTE list)}
\noindent[\textbf{New section.} Points that the fix note does not determine uniquely, or that are not settled; the text of the file was left as the author has it, with a red \textbf{NOTE} at the place.]
\begin{enumerate}
\item \textbf{DGLAP subtraction term} (Group II Row IV). fix\_poles\_dglap Step 11 writes $-\frac{2N_{c}}{\xi}\mathcal{F}(0)$; this equals the author's subtraction $-2N_{c}\frac{1-\xi}{\xi}\mathcal{F}(0)$ \emph{plus} the remnant ($-\frac{1}{k^{+}}\times$BFKL structure) of Group II Row IV, while Step 11 also keeps the remnant in Row IV. The author's convention (remnant in the finite terms of Row IV) is kept. Which writing is wanted for the paper?
\item \textbf{Fragmentation-function scheme} (Group II Row IV). \textbf{Answered} (author's answer (date 2026-11-20): ``Yes''.) The fragmentation functions are $\overline{\rm MS}$ number densities (pure-pole renormalisation with the $d$-dimensional Jacobian $\xi^{-d}$), i.e.\ the author's $\log\frac{(y'-y)^{2}\mu^{2}}{4\xi^{2}}$ form; this is now stated as the definition, and $\Delta$ to the CXY form, Eq.~\eqref{eq:fx-Delta}, is kept as a remark. (Not done, a cross-check only: an independent comparison with a standard NLO collinear $pp\to gX$ calculation.)
\item \textbf{Coordinate or local form of DGLAP and RC.} \textbf{Answered} (author's answer (date 2026-11-20): ``write it in both forms''.) Both forms are given, with their equivalence, Eqs.~\eqref{eq:fx-locsplit}, \eqref{eq:fx-equiv}: the coordinate forms are separately logarithmically IR divergent at large parent position, their sum with Group II Row I (second part) and Group III Row III is not (the tail cancels term by term, \texttt{notes/fix\_tail.tex}); in the local form the remainder $N_{\mathcal{R}}$ carries the tail.
\item \textbf{Order of limits.} The remaining pole $-\frac{\ell}{2}\frac{2}{\epsilon}$ vanishes only if $\vee\to\infty$ is taken before $\epsilon\to0$; this is assumed and should be stated in the paper.
\item \textbf{Two-$\epsilon$ versus single-$\epsilon$ bookkeeping.} Both give the same finite result; the paper should say which one is used and, if the single-$\epsilon$ rule is used, that the cancellation of the RC pole by the coupling counterterm rests on the scaleless normalisation integral (UV $=$ IR).
\end{enumerate}}\par
